\providecommand{\toplevelprefix}{../..}  %
\documentclass[../../book-main_es.tex]{subfiles}

\begin{document}

\chapter{Probemas pendientes y líneas de investigación en inteligencia}
\label{ch:future}

\begin{quote}
    ``\textit{El estudio se llevará a cabo partiendo de la hipótesis
        de que cada aspecto del aprendizaje o cualquier otra característica de
        la inteligencia puede describirse, en principio, con tal precisión que
        se pueda construir una máquina capaz de simularlo. Se intentará
        descubrir cómo hacer que las máquinas utilicen el lenguaje, formen abstracciones y
        conceptos, resuelvan tipos de problemas que actualmente están reservados a los seres humanos y
    se perfeccionen a sí mismas.}''

    $~$\hfill -- Proposición para el programa de IA de Dartmouth, 1956
\end{quote}
\vspace{5mm}

Este manuscrito presenta de manera sistemática los principios matemáticos y
los mecanismos computacionales que explican cómo se puede desarrollar la memoria o el 
conocimiento empírico a partir de datos observados de alta dimensión. La capacidad de 
buscar la parsimonia en un mundo aparentemente aleatorio es una característica fundamental 
de cualquier inteligencia, ya sea natural o artificial. Creemos que los principios y
mecanismos presentados en este libro son unificadores y universales, aplicables tanto 
a los animales como a las máquinas.

Esperamos que este libro ayude a los lectores a desentrañar por completo el misterio que rodea
las prácticas modernas de las redes neuronales profundas artificiales, mediante el desarrollo de una
comprensión rigurosa de sus funciones y su papel en el aprendizaje de (representaciones de)
distribuciones de baja dimensión a partir de datos de alta dimensión. Con esta
comprensión, quedan claras las capacidades y limitaciones de los modelos
y sistemas de IA existentes:
\begin{enumerate}
    \item Los modelos y sistemas de bucle abierto existentes no llegan a ser sistemas de memoria completos capaces de aprender y mejorar por sí mismos.
    \item Las implementaciones actuales del aprendizaje de representaciones siguen siendo rudimentarias y de fuerza bruta, lejos de ser óptimas en cuanto a arquitecturas de red o estrategias de optimización.
    \item Los modelos de IA actuales solo aprenden distribuciones de datos y
    realizan inferencias bayesianas inductivas, lo cual difiere de
    la inteligencia humana de alto nivel.
\end{enumerate}

Uno de los objetivos de este libro es ayudar a los lectores a adquirir una comprensión objetiva y
sistemática de las tecnologías actuales de inteligencia artificial
y a identificar los problemas y desafíos pendientes para seguir
avanzando en este campo. En el último capítulo, ofrecemos
algunas de nuestras opiniones y previsiones para el futuro.

\section{¿Hacia la inteligencia autónoma: cerrar el círculo?}
A partir de la práctica de la inteligencia artificial en la última década, ha
quedado claro que, si se cuenta con suficientes datos y recursos computacionales,
es posible construir un modelo lo suficientemente grande y preentrenarlo para que aprenda la
distribución \textit{a priori} de todos los datos, es decir, \(p(\vx)\).
En teoría, un modelo tan grande puede memorizar casi todo el
conocimiento existente sobre el mundo que está codificado en datos como imágenes 2D,
formas 3D, movimientos dinámicos y lenguajes naturales. Como comentamos al
principio del libro, un modelo tan amplio desempeña un papel similar al del ADN, que
la vida utiliza para registrar y transmitir el conocimiento sobre el mundo. 

El modelo y la distribución aprendidos de esta manera pueden utilizarse para
crear nuevas muestras de datos extraídas de la misma distribución. También se puede
utilizar el modelo para realizar inferencias (por ejemplo, completamiento,
estimación y predicción) con el conocimiento memorizado bajo
diversas condiciones, por ejemplo, mediante el muestreo de la distribución \textit{posterior}
\(p(\vx\mid \vy)\) bajo una nueva observación \(\vy\).
En sentido estricto, dicha inferencia es inductiva o estadística.

Ningún modelo preentrenado, por muy grande que sea, puede garantizar que la
distribución que ha aprendido hasta ahora sea totalmente correcta o completa.
Si nuestras muestras \(\hat{\vx}_{t}\) del \textit{a priori}
\(p_{t}(\vx)\) o las estimaciones \(\hat{\vx}_{t}(\vy)\) basadas en la
\textit{posterior} \(p_{t}(\vx\mid \vy)\) son inconsistentes con la
verdad \(\vx\), nos gustaría corregir las distribuciones aprendidas:
\begin{equation}
    p_{t}(\vx) \rightarrow p_{t+1}(\vx) \quad \text{or}\quad
    p_{t}(\vx\mid \vy) \rightarrow p_{t+1}(\vx\mid \vy),
\end{equation}
basado en el error \(\ve_{t} = \vx_{t} - \hat{\vx}_{t}\). Esto se
conoce como corrección de errores basada en la retroalimentación del error, un mecanismo
omnipresente en la naturaleza para el aprendizaje continuo. Sin embargo, cualquier modelo abierto
carece en sí mismo del mecanismo para revisar o mejorar la distribución aprendida
cuando esta es incorrecta o incompleta. La mejora de los modelos actuales de IA
sigue dependiendo en gran medida de la participación humana: supervisión o
refuerzo a través de la experimentación, la evaluación y la selección. Podríamos
denominar a este proceso «selección artificial» de modelos grandes, en
contraposición a la selección natural que rige la evolución de la vida.

Como hemos visto anteriormente en este libro (en
particular en el \Cref{ch:consistent-representations}), los sistemas de bucle cerrado
alinean una representación interna
con las observaciones del mundo externo. Pueden mejorar continuamente
la distribución aprendida internamente y su representación para
alcanzar la consistencia o la autoconsistencia. Un siguiente paso inmediato es
desarrollar sistemas de memoria de bucle cerrado, tal como se ilustra en
\Cref{fig:open-to-closed}, que aprendan de forma autónoma y continua
las distribuciones generales de los datos y mejoren a partir de la retroalimentación de errores.

Por lo tanto, la transición de los modelos de bucle abierto entrenados de extremo a extremo,
tan populares en la actualidad, a los sistemas de bucle cerrado con aprendizaje continuo
\begin{equation}
    \text{\textbf{open-ended} models} \; \Longrightarrow \;
    \text{\textbf{closed-loop} systems}
\end{equation}
es la clave para que las máquinas puedan emular de verdad cómo el cerebro animal aprende
y aplica el conocimiento en un mundo abierto. Creemos que
\begin{quote}
    \begin{center}
        \textit{open-ended models are for a closed world, however
        large; \\ closed-loop systems are for an open world, however small.}
    \end{center}
\end{quote}
Por lo tanto, la denominada «inteligencia artificial general» (AGI) nunca podrá
alcanzarse simplemente con un sistema capaz de memorizar todos los datos y conocimientos
existentes sobre el mundo. Ningún conocimiento, por mucho que sea, es verdaderamente generalizable.\footnote{Según Sir Karl Popper,
un influyente filósofo de la ciencia, ¡todas las teorías y conocimientos científicos son falsables!} Sin embargo,
es la capacidad de mejorar la memoria o el conocimiento existente y adaptarse a cualquier entorno nuevo lo que es
verdaderamente generalizable. Por lo tanto, la inteligencia verdaderamente genuina es en sí misma generalizable si se alcanza
de manera completa y correcta.\footnote{En nuestra opinión, la IGA debería representar la «inteligencia artificial genuina».}
Esta es precisamente la esencia del programa de cibernética esbozado por Norbert Wiener en la
década de 1940, del que hablamos al principio de este libro.

\begin{figure}[t]
    \centering
    \includegraphics[width=0.9\linewidth]{\toplevelprefix/chapters/chapter9/figs/open-to-closed.png}
    \caption{De una red profunda abierta a un sistema de circuito cerrado.}
    \label{fig:open-to-closed}
\end{figure}

\section{Hacia la inteligencia natural: ¿Más allá de la retropropagación?}
La práctica de la inteligencia artificial en los últimos años ha llevado a muchos a
creer que es necesario construir un único modelo de gran tamaño para aprender la
distribución de todos los datos y memorizar todo el conocimiento. Aunque esto
es tecnológicamente posible, es probable que tal solución esté lejos de ser
necesaria o eficiente. Como sabemos por el entrenamiento de redes profundas, el
único método escalable conocido para entrenar este tipo de redes a gran escala es a través de
la retropropagación (BP) \citep{Back-Prop}. Aunque la retropropagación (BP) ofrece una forma de
corregir errores mediante señales de gradiente que se propagan hacia atrás a través de todo el
modelo, se trata, no obstante, de un método bastante burdo y difiere
significativamente de la forma en que aprende la naturaleza: la BP es una opción que la naturaleza
no puede permitirse debido a su elevado costo y que, sencillamente, no puede implementar debido a
limitaciones físicas.

En términos más generales, no podemos comprender verdaderamente la inteligencia a menos que
también comprendamos cómo se puede implementar de manera eficiente. Es decir,
es necesario abordar la complejidad computacional que implica la implementación de los mecanismos
asociados al logro de los objetivos de la inteligencia. 
Históricamente, nuestra comprensión de la inteligencia (artificial) ha evolucionado
a través de varias fases, desde la complejidad de Kolmogorov —incomputable—
hasta la entropía de Shannon, desde la computabilidad de Turing hasta la posterior
comprensión de la tractabilidad,\footnote{Decimos que un problema es tractable
    si permite un algoritmo cuya complejidad es polinómica en el tamaño
del problema.} y hasta el fuerte énfasis en la escalabilidad de los algoritmos
en la práctica moderna del análisis de datos de alta dimensión
\citep{Wright-Ma-2022} y la inteligencia artificial. Esta evolución se puede
resumir de la siguiente manera:
\begin{equation}
    \text{\textbf{incomputable}} \;
    \Longrightarrow \; \text{\textbf{computable}} \;
    \Longrightarrow \; \text{\textbf{tractable}} \; \Longrightarrow \;
    \text{\textbf{scalable}}.
\end{equation}
En gran medida, el éxito y la popularidad del aprendizaje profundo y
la retropropagación se deben precisamente a que han ofrecido una implementación
escalable con plataformas informáticas modernas (como las GPU) para
procesar y comprimir grandes volúmenes de datos. Sin embargo, dicha
implementación sigue siendo mucho más costosa que la forma en que la naturaleza
llega a la inteligencia.

Aún queda un margen considerable para mejorar la eficiencia de la inteligencia artificial,
de modo que pueda emular la eficiencia de la inteligencia natural,
la cual debería ser varios órdenes de magnitud superior a
las actuales implementaciones basadas en la fuerza bruta. Para ello, necesitamos descubrir
nuevas arquitecturas de aprendizaje y mecanismos de optimización que permitan
aprender distribuciones de datos en condiciones físicas naturales y
bajo restricciones de recursos, similares a las que enfrentan los seres inteligentes en
la naturaleza; por ejemplo, sin acceder a todos los datos de una sola vez ni actualizar
todos los parámetros del modelo simultáneamente (mediante BP).

El marco teórico y el enfoque que se exponen en este libro pueden servirnos de guía
para descubrir esas nuevas arquitecturas y mecanismos. Estas nuevas
arquitecturas y mecanismos deberían permitir el aprendizaje continuo en línea
y deberían poder actualizarse mediante optimización hacia adelante
o hacia atrás altamente localizada y dispersa.\footnote{Para aprender una distribución, una
    aplicación sencilla de estos requisitos se da en el caso más simple del
    PCA, con el método de PCA en línea presentado en el
    \Cref{ch:consistent-representations}. Una vez aprendido el modelo lineal de la
    distribución, la tarea más sencilla de realizar predicciones en línea
    con observaciones ruidosas puede llevarse a cabo fácilmente mediante métodos de mínimos cuadrados
    , como el filtro de Wiener \cite{Wiener-1942,Wiener-1949}. En el caso de los sistemas dinámicos lineales, esto se convierte en el famoso
filtro de Kalman; véase \cite[Apéndice B]{MaY2003} o \citep{peng2025mathematics}.}
  
\begin{figure}[t]
    \centering
    \includegraphics[width=0.99\linewidth]{\toplevelprefix/chapters/chapter9/figs/loops.png}
    \caption{Estructura hipotética de la corteza cerebral. La corteza
        es un sistema de autoencodificación masivamente paralelo y distribuido
        que consiste en una jerarquía de autoencodificadores de bucle cerrado. Estos
        extraen información de sus entradas y maximizan la
        ganancia de información de las representaciones que generan. Dichos procesos se llevan a cabo en múltiples
    niveles de la jerarquía con distintos grados de abstracción de la información.}
    \label{fig:loops}
\end{figure}
Tal y como nos ha enseñado la neurociencia, la corteza de nuestro cerebro
está formada por decenas de miles de columnas corticales
\citep{Hawkins-2021}. Todas las columnas corticales tienen estructuras
físicas y funciones similares. Son altamente paralelas y están distribuidas,
aunque están interconectadas de forma dispersa. Por lo tanto, creemos que para desarrollar un
sistema de memoria más escalable y estructurado, debemos considerar
arquitecturas que emulen las de la corteza. La \Cref{fig:loops}
muestra una arquitectura hipotética de este tipo: un sistema masivamente distribuido y
jerárquico compuesto por numerosos módulos de autoencodificación de bucle cerrado
que funcionan en gran medida en paralelo.\footnote{En la bibliografía existen elementos de este tipo de
    arquitecturas hipotéticas en diversas formas, como
    las redes de codificación predictiva \citep{millidge2022predictive}, que cuentan
con una estrategia de optimización que utiliza actualizaciones locales eficientes.}
Estos módulos aprenden a codificar diferentes modalidades sensoriales o múltiples
proyecciones de datos de cada modalidad sensorial. Nuestro análisis en
\Cref{sec:measurement-self-consistency} de
\Cref{ch:conditional-inference} sugiere que dicha detección
y aprendizaje paralelos de una distribución de baja dimensión son teóricamente
posibles. Higher-level (lossy) autoencoders can then be learned based
on outputs of lower-level ones to develop more sparse and
higher-level ``abstractions'' of the representations learned by the
lower levels.

La arquitectura de sistema distribuida, jerárquica y de bucle cerrado
que se ilustra en la \Cref{fig:loops} comparte muchas características con la
corteza cerebral. Dicha arquitectura de sistema puede abrir muchas más
posibilidades que la actual arquitectura de un único modelo grande. Permite
explorar mecanismos de aprendizaje y optimización mucho más eficientes
y da como resultado una organización modular más estructurada de
la distribución de los datos aprendidos y el conocimiento, emulando más fielmente la memoria
de nuestro cerebro. Esto nos permitiría llevar la implementación de la inteligencia artificial
al siguiente nivel de evolución:
\begin{equation}
    \text{\textbf{incomputable}}
    \Longrightarrow \text{\textbf{computable}}
    \Longrightarrow \text{\textbf{tractable}} \Longrightarrow
    \text{\textbf{scalable}}
    \Longrightarrow
    \text{\textbf{natural}}.
\end{equation}

\section{¿Hacia la inteligencia científica: más allá del test de Turing?}

\subsection{La evolución de la inteligencia en la naturaleza}
Como hemos comentado al principio de este libro, en el \Cref{ch:intro},
la inteligencia en la naturaleza ha evolucionado a lo largo de múltiples fases y
se ha manifestado en cuatro formas distintas:
\begin{equation}
    \text{\textbf{phylogenetic}} \;
    \Longrightarrow \; \text{\textbf{ontogenetic}} \;
    \Longrightarrow \; \text{\textbf{societal}} \;
    \Longrightarrow \; \text{\textbf{scientific intelligence}}.
\end{equation}
Todas las formas de inteligencia comparten el objetivo (y la
característica) común de adquirir conocimientos útiles en forma de distribuciones de baja dimensión
a partir de datos de alta dimensión sobre el mundo que se perciben.
No obstante, difieren en los esquemas de codificación específicos que adoptan, la
información codificada, los mecanismos computacionales para el aprendizaje y
la mejora, y las implementaciones físicas de dichos mecanismos. Según
los conceptos y la terminología desarrollados en este libro, las cuatro etapas
de la inteligencia se diferencian en los tres aspectos siguientes:
\begin{enumerate}
    \item El \textit{libro de códigos} o esquema utilizado para codificar la
        información o el conocimiento que se desea transmitir.
    \item La \textit{información} o los conocimientos adquiridos y codificados
        mediante el libro de códigos.
    \item Los \textit{mecanismos de optimización} utilizados para crear y
        corregir la información o el conocimiento codificados.
\end{enumerate}
La \Cref{tab:four-stage-intelligence} resume sus principales características.
\begin{table}[h]
    \centering
    \resizebox{1\columnwidth}{!}
    {
        \begin{tabular}{| c | c | c | c | c |}
            \hline & \textbf{Filogenéico} & \textbf{Ontogenético} &
            \textbf{Social} & \textbf{Científico}\\
            \hline

            \textbf{Libro de código}  & DNAs & Neuronas/Mente & Lenguajes Naturales
            & Abstracciones matemáticas \\ [0.5ex]
            \hline

            \textbf{Información} & Genes & Memoris & Conocimiento
            empírico & Factos científicos\\ [0.5ex]
            \hline
            \textbf{Optimización} & Refuerzo & Retroalimentación de error &
            Prueba \& error & Teorizar \& Falsar \\  [0.5ex]
            \hline
        \end{tabular}
    }
    \caption{Características principales de las cuatro etapas de la inteligencia en la naturaleza.}
    \label{tab:four-stage-intelligence}
\end{table}

Como ya sabemos, los seres humanos han logrado dos avances decisivos en materia de inteligencia.
\begin{enumerate}
    \item El primero fue el desarrollo del lenguaje oral y escrito,
        lo que permitió a los seres humanos compartir y transmitir el conocimiento adquirido
        de una generación a otra, de manera muy similar a como lo hace el ADN en la naturaleza.
    \item El segundo fue el desarrollo de las matemáticas y la lógica formal
        hace aproximadamente tres mil años, que se convirtieron en el
        lenguaje preciso de la ciencia moderna.
\end{enumerate}
El lenguaje de las matemáticas nos liberó de resumir el conocimiento a partir de
las observaciones únicamente de forma empírica y nos permitió formalizar
el conocimiento en forma de teorías verificables o falsificables mediante la deducción
matemática o la experimentación. A través de la formulación de hipótesis, la deducción
lógica y la verificación o falsificación experimental, ahora podemos
descubrir y desarrollar de manera proactiva nuevos conocimientos que van mucho más allá
de lo que se puede aprender a partir de las distribuciones de los datos
observados.\footnote{Por ejemplo, las relaciones causales no pueden deducirse
únicamente a partir de las distribuciones \citep{Pearl-2009}.}

A pesar de las aparentes diferencias entre estas distintas etapas de
la inteligencia, la evolución de la inteligencia comparte una característica
común:\textit{Todos ellos utilizan ciertas formas de
    retroalimentación\footnote{Pérdida por error y recompensa por éxito.} para crear
    y corregir la información aprendida, y la retroalimentación y la corrección son
    cada vez más frecuentes y eficientes a lo largo de la evolución: desde el refuerzo de bucle abierto hasta la autosupervisión de bucle cerrado, desde
el ensayo y error colectivo hasta la formulación proactiva de hipótesis y su refutación.}

\subsection{De la inteligencia inductiva a la deductiva}
Como se mencionó en la introducción (\Cref{ch:intro}), el programa de
«inteligencia artificial» (IA) de 1956 tenía precisamente como objetivo estudiar cómo
dotar a las máquinas de inteligencia científica, es decir, de funciones de alto nivel
como la abstracción matemática, la inferencia lógica y
la resolución de problemas, que se cree que diferencian a los seres humanos de los animales:
\begin{equation}
    \text{\textbf{low-level} (animal) intelligence} \; \Longrightarrow \;
    \text{\textbf{high-level} (human) intelligence.}
\end{equation}
Como hemos aclarado en repetidas ocasiones en este libro, la mayoría de los avances tecnológicos
en inteligencia artificial durante la última década, aunque
se hayan llevado a cabo bajo el nombre de «IA», están en realidad más estrechamente relacionados
con el hecho de que las máquinas imiten formas de inteligencia de bajo nivel que comparten en gran medida
tanto los animales como los seres humanos, incluyendo la inteligencia filogenética,
ontogenética y social. En estos niveles, la inteligencia
crea memoria o conocimiento \textit{empírico} (a partir de distribuciones de datos empíricos) y lleva a cabo inferencias
que son, en su mayoría, de naturaleza \textit{inductiva}. Los principios y métodos presentados en este libro tienen como objetivo revelar los objetivos científicos y
los mecanismos computacionales que subyacen a la inteligencia en estos niveles. Estos
proporcionan una sólida justificación teórica y evidencia computacional
para la afirmación:
\begin{quote}
    \centering
    \textit{Las máquinas pueden adquirir conocimientos empíricos y realizar
    inferencias inductivas.}
\end{quote}
Sin embargo, hasta ahora no hay pruebas teóricas ni científicas rigurosas
que indiquen que estos mecanismos por sí solos bastarían para alcanzar
el alto nivel de inteligencia humana que el programa original de IA realmente
pretendía comprender y quería que las máquinas imitaran o incluso superaran.

De hecho, a día de hoy, sabemos muy poco sobre cómo verificar
o certificar de manera rigurosa si un sistema es realmente capaz de alcanzar un nivel tan alto
de inteligencia, a pesar de que el Test de Turing, también conocido como el
Juego de la Imitación,\footnote{En la propuesta de Turing, el evaluador, o
    interrogador, es un ser humano. Sin embargo, la mayoría de los evaluadores humanos tienen una
    formación y unos conocimientos científicos limitados, y sus conclusiones pueden ser
subjetivas.} se propuso en 1950 \citep{Turing-1950}, como un primer
intento de abordar la pregunta fundamental:
\begin{quote}
    \centering
    \textit{¿Pueden pensar las máquinas?}
\end{quote}
Durante mucho tiempo, no se consideró
necesario establecer una definición precisa de dicha prueba, ya que las capacidades de las máquinas estaban muy por debajo de las de los seres humanos
(o incluso de las de los animales). Sin embargo, dados los recientes avances tecnológicos, muchos
modelos y sistemas afirman ahora alcanzar o incluso superar la inteligencia
humana. Por lo tanto, ya es hora de desarrollar una definición científica y
aplicable de la prueba de Turing, es decir, un protocolo sistemático y
riguroso para evaluar y certificar el nivel de inteligencia de un
modelo o sistema que afirme ser «inteligente».

Como argumentó Turing en su artículo original, resulta bastante difícil
definir qué es «pensar». En su lugar, podríamos partir de algo más
tangible e intentar determinar si un sistema es realmente capaz de
«comprender» ciertos conceptos o conocimientos. Por ejemplo, ¿cómo podemos
verificar rigurosamente si un modelo o sistema inteligente ha comprendido realmente
un concepto abstracto como: 
\begin{enumerate}
    \item el concepto de igualdad entre dos cantidades,
    \item el concepto de los números (naturales, racionales o reales),
    \item el concepto de infinito,
    \item y el concepto de inducción matemática,
\end{enumerate}
por mencionar solo algunos? ¿O simplemente ha memorizado una enorme cantidad de
ejemplos de tales conceptos? Cabe señalar que los modelos de lenguaje de última generación
siguen teniendo dificultades con preguntas matemáticas sencillas como: «¿Es
3,11 mayor o menor que 3,9?»\footnote{Algunos modelos han
    corregido sus respuestas a este tipo de preguntas mediante
    ingeniería específica, o han incorporado mecanismos de razonamiento adicionales
    que verifican las respuestas inmediatas y las corrigen durante la inferencia.
    Sin embargo, dejamos al lector la tarea de comprobar rigurosamente
    si algún modelo lingüístico de vanguardia comprende realmente el
    concepto de los números (naturales, racionales, reales y complejos) y su
    aritmética asociada.}

Gracias a este libro, sabemos que los sistemas modernos de IA generativa se basan principalmente
en la inferencia bayesiana (y, por lo tanto, inductiva) para generar sus respuestas. ¿Pueden
las máquinas también comprender verdaderamente la inducción matemática o la deducción lógica
general y entender su diferencia con respecto a la inducción bayesiana?
¿Cómo verificamos si un sistema comprende realmente las reglas de
la deducción lógica y matemática y es capaz de aplicarlas con rigor? ¿O,
una vez más, simplemente ha memorizado una gran cantidad de ejemplos de tales
deducciones —como datos sobre la cadena de pensamiento— y ha aprendido y
aprovechado sus patrones de distribución? Por lo tanto, nos queda una pregunta sin respuesta:
\begin{quote}
    \centering
    \textit{¿Hay alguna diferencia entre \textbf{memorizar} y
    \textbf{comprender}?}
\end{quote}
Solo conociendo la diferencia podremos intentar dar una
respuesta significativa a la pregunta: \textit{«¿Pueden las máquinas entender?»}

Probablemente, lo más importante sea preguntarnos: ¿cómo verificamos si un sistema que se autodenomina
inteligente es capaz de \textit{crear conceptos abstractos o mejorar su propio conocimiento} de manera proactiva mediante métodos deductivos? Por ejemplo, ¿puede crear nuevos conceptos matemáticos
y descubrir nuevas leyes físicas o relaciones causales?
La conexión entre comprender y crear puede ser bastante fundamental, como lo expresó de manera célebre el físico Richard Feynman:
\begin{quote}
    \centering
    \textit{«Lo que no puedo crear, no lo entiendo».}
\end{quote}

\subsection{Pruebas científicas de inteligencia}
Por lo tanto, ya es hora de que desarrollemos métodos de evaluación con base científica
para determinar a cuál de las siguientes categorías pertenece la capacidad
inteligente de un sistema:
\begin{enumerate}[label={Level-\arabic*.}, leftmargin=*]
    \item el mero hecho de haber aprendido la distribución de ciertos
        datos que contienen información y ser capaz de reproducirlos;
    \item ser capaz de aprender de forma autónoma y continua
        nuevos conocimientos o de corregir los conocimientos existentes a partir de nuevas
        observaciones y experiencias;
    \item ser capaz de comprender y crear verdaderamente
        conceptos, reglas y leyes abstractas, y saber cómo aplicarlos
        correctamente de manera deductiva;
    \item ser capaz de formular nuevas hipótesis científicas o conjeturas matemáticas y verificarlas o refutarlas mediante la experimentación o la deducción lógica.
\end{enumerate}

\begin{figure}[t]
    \centering
    \includegraphics[width=\linewidth]{\toplevelprefix/chapters/chapter9/figs/tests.png}
    \caption{Tres pruebas para diferentes niveles o tipos de
        capacidades intelectuales: la prueba de Wiener para la inteligencia básica,
        la prueba de Turing para la inteligencia de nivel humano y
    la prueba de Popper para la inteligencia de nivel científico.}
    \label{fig:three-tests}
\end{figure}

Para evaluar y distinguir estos diferentes niveles de inteligencia,
sugerimos que probablemente deberían existir al menos tres pruebas diferentes,
tal y como se ilustra en la \Cref{fig:three-tests}:
\begin{enumerate}
    \item \textit{La prueba de Norbert Wiener:} para determinar si un
        sistema es capaz de autocorregirse y generar nuevo conocimiento (empírico)
        por sí mismo o si simplemente recibe y actualiza información
        de manera pasiva a través de refuerzos o supervisión externos;
    \item \textit{El test de Alan Turing:} para determinar si un
        sistema es capaz de crear y comprender conceptos abstractos y
        conocimiento, o si simplemente aprende y memoriza estadísticas de las
        muestras y las utiliza para la inferencia bayesiana;
    \item \textit{La prueba de Karl Popper:} para determinar si un
        sistema es capaz de explorar y crear de forma proactiva nuevo conocimiento abstracto
        mediante la formulación y verificación de hipótesis o teorías
        basadas en la consistencia experimental o la coherencia lógica interna.
\end{enumerate}
Creemos que, para este tipo de pruebas, el evaluador o el interrogador
no debería ser una sola persona ni un grupo arbitrario de personas, sino más bien una
entidad con acreditación científica que siga un protocolo y un proceso
científicamente sólidos, de modo que las conclusiones sean rigurosas y fiables.

Como hemos visto a lo largo de este libro, la \textit{compresión} ha desempeñado
un papel fundamental en el desarrollo de la memoria o en la creación de conocimiento 
empírico. Es el principio rector y un mecanismo universal para identificar una
distribución de datos empíricos y organizar la información codificada en ellos 
mediante una representación estructurada. En gran medida, explica la práctica de la 
«inteligencia artificial» con redes profundas, que es en gran parte de naturaleza 
inductiva.\footnote{Los sistemas que cuentan con mecanismos deductivos integrados para 
llevar a cabo la verificación formal no entran en esta categoría.} Una cuestión pendiente para futuros
estudios es si \textit{la compresión por sí sola} es suficiente para alcanzar
toda la inteligencia deductiva de alto nivel mencionada anteriormente. 
En concreto, ¿son la abstracción matemática, la inferencia causal, la generación de hipótesis
y el razonamiento y la deducción lógicos algún tipo de formas ampliadas o
superadas de la compresión? ¿Existe alguna diferencia fundamental entre aprender distribuciones de datos 
a través de la compresión y formar conceptos y teorías de alto nivel mediante la abstracción? De ahí que 
tengamos otra pregunta sin respuesta:
\begin{quote}
    \centering
    \textit{¿Hay alguna diferencia entre la \textbf{compresión} y
    la \textbf{abstracción}?}
\end{quote}
Solo conociendo la diferencia podremos intentar dar una respuesta significativa a la pregunta: 
«¿Pueden las máquinas crear conocimiento de forma proactiva mediante la deducción?»

En gran medida, la ciencia---y su libro de códigos asociado,
las matemáticas---puede considerarse como la forma más avanzada de
inteligencia, exclusiva de los seres humanos cultos e ilustrados. 
El filósofo Sir Karl Popper sugirió en una ocasión:
\begin{quote}
    \centering
    \textit{«La ciencia puede describirse como el arte de la
    simplificación excesiva.»}
\end{quote}
Creemos que descubrir y comprender los principios matemáticos y los mecanismos computacionales subyacentes
de esa inteligencia de alto nivel, y reproducirlos con éxito mediante máquinas,
será la última frontera para la ciencia, las matemáticas y la computación en su conjunto.

\end{document}
